\section{工业化与尺度：frontier model 为什么不能直接复现}

本节从课程方法论转向现实约束：语言模型已经工业化。GPT-4 级别模型训练成本据称达到上亿美元；xAI 等公司建设数十万 GPU 集群。这些事实改变了教学策略：课堂不能靠复现 frontier run 学习，而要靠小规模实验、资源核算和公开模型证据训练判断力。

\begin{knowledgebox}{课堂提示：小模型实验要明确哪些结论可外推}
老师强调，小模型最适合验证 mechanics、资源公式和相对趋势；绝对最优超参、通信占比、能力阈值与数据收益可能随 scale 改变。每个实验都应写清“机制结论”和“尺度假设”，避免把课堂可复现实验包装成 frontier recipe。
\end{knowledgebox}

\begin{figure}[H]
\centering
\includegraphics[width=0.74\textwidth]{industrialisation.jpg}
\caption{工业化隐喻：前沿语言模型训练已经是大型基础设施工程。}
\figsource{CS336 Spring 2026 Lecture 1 官方可执行讲义，\texttt{why\_this\_course\_exists()}，原图来自 Wikimedia Commons。}
\end{figure}

\begin{knowledgebox}{读图：工业化图到底说明什么}
图中是传统工业生产场景。它在这里对应 LLM 的集群、电力、网络、数据、冷却、工程团队和资本投入。模型质量不再只由一个算法公式决定，而是由整套生产系统决定。后续 systems、scaling laws、data 单元都在拆这个生产系统。
\end{knowledgebox}

闭源 frontier model 的另一个教学障碍是细节不可见。GPT-4 technical report 明确隐藏了模型大小、硬件、训练 compute、数据构成和训练方法等关键细节。

\begin{figure}[H]
\centering
\includegraphics[width=0.84\textwidth]{gpt4-no-details.png}
\caption{GPT-4 technical report 对训练细节的隐藏，说明 frontier model 已无法作为完整可复现实验对象。}
\figsource{CS336 Spring 2026 Lecture 1 官方可执行讲义，\texttt{why\_this\_course\_exists()}。}
\end{figure}

\begin{warningbox}{读图时容易犯的错误}
这张图不是在讨论开放/闭源的价值判断，而是在说明课程约束：如果 frontier recipe 不透明，学生就不能通过照抄 recipe 学到完整技术栈。更可靠的学习目标是 mechanics、resource accounting、scaling mindset 和实验设计。
\end{warningbox}

\subsection{尺度会改变系统瓶颈}

前面的工业化约束首先体现在系统账本：课堂能训练小于 1B 参数的模型，但这不代表小模型结论能直接外推。随着模型变大，FLOPs 分布、显存瓶颈、通信瓶颈、batch shape、推理成本都会变；因此后续 systems 单元会同时做静态资源核算和真实 runtime profiling。

\begin{figure}[H]
\centering
\includegraphics[width=0.72\textwidth]{roller-flops.png}
\caption{不同模型尺度下 attention、FFN、logit 等模块的 FLOPs 占比会变化。}
\figsource{CS336 Spring 2026 Lecture 1 官方可执行讲义，图片 \texttt{images/roller-flops.png}。}
\end{figure}

\begin{knowledgebox}{读图：Roller FLOPs 表应该怎么看}
每一行是一个 OPT 模型尺度，例如 760M、1.3B、13B、175B。列中 MHA 表示 multi-head attention 成本，FFN 表示 feed-forward network 成本，attn 更聚焦 attention score/value mixing，logit 是输出词表投影。关键趋势是：模型越大，FFN 占比越高，logit 和 attention 的相对占比下降。它说明“Transformer”不是固定资源画像；同一架构名下的真实瓶颈会随尺度漂移。
\end{knowledgebox}

\begin{warningbox}{FLOPs 占比不是 wall-clock 时间占比}
FLOPs 少的模块仍可能因为 memory bandwidth、kernel launch、通信、KV cache 读写而拖慢系统。读这张图时要区分 compute accounting 和 runtime profiling。Lecture 2 会建立这两者之间的桥。
\end{warningbox}

\subsection{尺度也会改变能力表现}

系统瓶颈之外，本节进一步看能力曲线。讲义引用 emergence 相关图：许多任务在小规模区域贴近 random baseline，跨过某个训练规模后指标突然上升。它的教学信号不是“大模型有魔法”，而是小规模实验看到的能力曲线可能不完整，评测分辨率和指标形式也可能制造表面突变。

\begin{knowledgebox}{课堂提示：不要把一张 emergence 图当成因果解释}
课上提醒，这类图只显示 metric 随 scale 的观测关系。要判断能力是否真的突然出现，还需检查连续 loss、评分阈值、prompt 格式和模型族差异；讲义保留它是为了训练外推谨慎，而不是宣布所有能力都存在神秘临界点。
\end{knowledgebox}

\begin{figure}[H]
\centering
\includegraphics[width=0.84\textwidth]{wei-emergence-plot.png}
\caption{一些任务指标随 training FLOPs 出现陡峭上升，提示小模型直觉不能无脑外推。}
\figsource{CS336 Spring 2026 Lecture 1 官方可执行讲义，图片 \texttt{images/wei-emergence-plot.png}。}
\end{figure}

\begin{knowledgebox}{读图：emergence 曲线应该怎么看}
每个小图是一项任务，横轴是 training FLOPs 且为 log scale，纵轴是 accuracy、BLEU 或 exact match。不同颜色是不同模型族，粉色虚线是 random baseline。许多曲线先长期贴着 baseline，然后在某个规模附近明显上升。这提醒我们：评测指标可能像阈值函数，模型能力也可能在小规模实验中不可见。
\end{knowledgebox}

\begin{warningbox}{emergence 图不该被过度解读}
陡峭上升不一定证明模型内部突然长出全新模块。它可能受指标离散化、prompt 格式、样本难度、横轴取对数、模型族差异影响。对 CS336 来说，真正结论是：mechanics 和 mindset 可迁移，intuition 需要谨慎外推。
\end{warningbox}

\subsection{三类可迁移知识}

第一讲把知识分成 mechanics、mindset、intuitions：

\begin{longtable}{p{0.20\textwidth}p{0.43\textwidth}p{0.27\textwidth}}
\toprule
\textbf{类型} & \textbf{含义} & \textbf{迁移性} \\
\midrule
Mechanics & Transformer、BPE、model parallelism、optimizer state、KV cache 等机制如何工作 & 高，可以在课堂规模实现和检查。 \\
Mindset & 如何做 resource accounting、如何压榨硬件、如何认真对待 scaling & 高，是跨模型版本的工程习惯。 \\
Intuitions & 哪些数据、架构、训练技巧通常有效 & 部分迁移，依赖尺度、数据和系统上下文。 \\
\bottomrule
\end{longtable}

\begin{figure}[H]
\centering
\includegraphics[width=0.82\textwidth]{divine-benevolence.png}
\caption{Noam Shazeer 对 SwiGLU 的表述提醒我们：一些成功设计首先来自实验，而不是事前完备理论。}
\figsource{CS336 Spring 2026 Lecture 1 官方可执行讲义，\texttt{why\_this\_course\_exists()}。}
\end{figure}

\begin{knowledgebox}{读图：这句玩笑背后的方法论}
SwiGLU 是后来许多 LLM 采用的 MLP 激活变体。图中文字的幽默点是：讲者承认某些设计最初不是由漂亮理论推出，而是由实验发现。课程要教的不是神秘化 intuition，而是如何设计实验、解释结果，并判断结论能否迁移。
\end{knowledgebox}

\subsection{Bitter Lesson 的正确读法}

课程把 The Bitter Lesson 的误读拆开：错误读法是“规模就是一切，算法不重要”；正确读法是“能随着资源增长持续受益的算法最重要”。

\[
\text{accuracy} = \text{efficiency} \times \text{resources}.
\]

\begin{itemize}
    \item \(\text{accuracy}\)：模型效果，可用 loss、held-out likelihood、benchmark 或真实任务质量衡量。
    \item \(\text{efficiency}\)：单位资源转化为有效能力的效率，包含算法、系统、数据和调参效率。
    \item \(\text{resources}\)：数据、训练 FLOPs、显存、通信带宽、工程时间和推理预算。
\end{itemize}

\begin{importantbox}{Bitter Lesson 对 CS336 的意义}
规模放大了效率的重要性。课堂小模型浪费一点算力只是慢几分钟；frontier-scale 训练中同样比例的浪费可能意味着数百万美元、数周排队和错过实验窗口。
\end{importantbox}

\subsection{本章小结}

工业化和尺度共同决定了 CS336 的路线：课堂不能复现 frontier model，但能训练可迁移的实现能力、资源意识和实验判断。小模型是学习机制的实验台，不是前沿模型的完美缩小版。

